\chapter{Elementary Functions}

\section{Polynomials}

\begin{definition}{Polynomial Function}{polynomial_def}
A \textbf{polynomial} of the complex variable $z$ is a function of the form
\begin{equation}
p_n(z) = a_0 + a_1 z + a_2 z^2 + \dots + a_n z^n = \sum_{k=0}^n a_k z^k,
\end{equation}
where $a_0, a_1, \dots, a_n \in \mathbb{C}$ are complex constants.
If $a_n \neq 0$, $n$ is called the \textbf{degree} of the polynomial.
Every polynomial is an \textbf{entire function}.
\end{definition}

\begin{theorem}{Factorization of Polynomials}{poly_factorization}
Every polynomial $p_n(z)$ of degree $n \ge 1$ can be factored completely into $n$ linear factors over $\mathbb{C}$:
\begin{equation}
p_n(z) = a_n(z - z_1)(z - z_2)\cdots(z - z_n),
\end{equation}
where $z_1, z_2, \dots, z_n \in \mathbb{C}$ are the zeros of $p_n(z)$ (counting multiplicity).
\begin{itemize}
    \item A polynomial of degree $n$ has exactly $n$ zeros (counting multiplicity).
    \item A polynomial is completely determined by its zeros up to the constant leading coefficient $a_n$.
    \item A point $z_0$ is a zero of multiplicity $k$ if and only if
    \begin{equation}
    p_n(z) = (z - z_0)^k q(z) \quad \text{with } q(z_0) \neq 0.
    \end{equation}
\end{itemize}
\end{theorem}

\begin{example}{Polynomial Factorization}{factor_ex}
\begin{enumerate}
    \item $p_3(z) = -2z^3 - 4z^2 + 10z + 12 = -2(z - 2)(z + 1)(z + 3)$. (All zeros are real: $2, -1, -3$).
    \item $p(z) = z^3 + (2 - i)z^2 - 2iz = z(z^2 + (2-i)z - 2i) = z(z + 2)(z - i)$. (Zeros: $0, -2, i$).
\end{enumerate}
\end{example}

\begin{proposition}{Polynomials with Real Coefficients}{real_poly_factors}
Let $p(z)$ be a polynomial with all coefficients real ($a_k \in \mathbb{R}$). If $z_0$ is a zero of $p(z)$, then its complex conjugate $\bar{z}_0$ is also a zero of $p(z)$ with the same multiplicity.
Consequently, any polynomial with real coefficients can be factored over $\mathbb{R}$ into a product of real linear factors and irreducible real quadratic factors of the form:
\begin{equation}
(z - z_1)(z - \bar{z}_1) = z^2 - 2\Re(z_1)z + |z_1|^2.
\end{equation}
\end{proposition}

\begin{example}{Real Factorization of $z^5 - 1$}{z5_factor}
The fifth roots of unity yield:
\begin{equation}
z^5 - 1 = (z - 1)\left(z^2 - 2\cos\left(\frac{2\pi}{5}\right)z + 1\right)\left(z^2 - 2\cos\left(\frac{4\pi}{5}\right)z + 1\right).
\end{equation}
\end{example}

\begin{theorem}{Taylor Expansion of a Polynomial}{poly_taylor}
Every polynomial $p_n(z)$ of degree $n$ can be expanded in powers of $(z - z_0)$ about any point $z_0 \in \mathbb{C}$:
\begin{equation}
p_n(z) = \sum_{k=0}^n \frac{p_n^{(k)}(z_0)}{k!} (z - z_0)^k.
\end{equation}
\begin{itemize}
    \item If $z_0 = 0$, this is the \textbf{Maclaurin expansion}.
    \item $z_0$ is a zero of $p_n(z)$ if and only if $p_n(z_0) = 0$.
    \item $z_0$ is a zero of \textbf{multiplicity $k$} if and only if:
    \begin{equation}
    p_n(z_0) = p'_n(z_0) = \dots = p_n^{(k-1)}(z_0) = 0 \quad \text{and} \quad p_n^{(k)}(z_0) \neq 0.
    \end{equation}
\end{itemize}
\end{theorem}


\section{Rational Functions and Partial Fraction Decomposition}

\subsection{General Rational Functions}

\begin{definition}{Rational Function, Zeros, and Poles}{rational_func_def}
A \textbf{rational function} is the ratio of two polynomials:
\begin{equation}
R_{m, n}(z) = \frac{P_m(z)}{Q_n(z)} = \frac{a_0 + a_1 z + \dots + a_m z^m}{b_0 + b_1 z + \dots + b_n z^n} \quad (a_m \neq 0, \, b_n \neq 0).
\end{equation}
Assuming all common factors have been canceled:
\begin{itemize}
    \item The zeros of $P_m(z)$ are the \textbf{zeros} of $R_{m, n}(z)$.
    \item The zeros of $Q_n(z)$ are the \textbf{poles} of $R_{m, n}(z)$ (points where $R_{m, n}(z) \to \infty$).
    \item $R_{m, n}(z)$ is analytic everywhere in $\mathbb{C}$ except at its poles.
\end{itemize}
\end{definition}

\begin{proposition}{Polynomial Division for Rational Functions}{rat_poly_div}
If the degree of the numerator $m$ is greater than or equal to the degree of the denominator $n$ ($m \ge n$), Euclidean polynomial division guarantees the unique decomposition:
\begin{equation}
\frac{P_m(z)}{Q_n(z)} = E_{m-n}(z) + \frac{F_\ell(z)}{Q_n(z)},
\end{equation}
where $E_{m-n}(z)$ is a polynomial of degree $m-n$, and $F_\ell(z)$ is a polynomial of degree $\ell < n$.
\end{proposition}

\subsection{Partial Fraction Decomposition}

\begin{theorem}{Partial Fraction Decomposition}{pfd_theorem}
Let $R(z) = \frac{P(z)}{Q(z)}$ be a strictly proper rational function ($\deg P < \deg Q = n$). Let the distinct poles of $R(z)$ be $\zeta_1, \zeta_2, \dots, \zeta_r$ with respective multiplicities $d_1, d_2, \dots, d_r$ (such that $\sum_{j=1}^r d_j = n$):
\begin{equation}
Q(z) = b_n (z - \zeta_1)^{d_1} (z - \zeta_2)^{d_2} \cdots (z - \zeta_r)^{d_r}.
\end{equation}
Then $R(z)$ has a unique \textbf{partial fraction decomposition}:
\begin{equation}
R(z) = \sum_{j=1}^r \left[ \frac{A_0^{(j)}}{(z - \zeta_j)^{d_j}} + \frac{A_1^{(j)}}{(z - \zeta_j)^{d_j - 1}} + \dots + \frac{A_{d_j - 1}^{(j)}}{z - \zeta_j} \right] = \sum_{j=1}^r \sum_{s=0}^{d_j - 1} \frac{A_s^{(j)}}{(z - \zeta_j)^{d_j - s}},
\end{equation}
where the coefficients are given by the \textbf{derivative method}:
\begin{equation}
A_s^{(j)} = \lim_{z \to \zeta_j} \frac{1}{s!} \frac{d^s}{dz^s} \left[ (z - \zeta_j)^{d_j} R(z) \right], \quad s = 0, 1, \dots, d_j - 1.
\end{equation}
\end{theorem}

\begin{example}{Partial Fraction Decomposition with Repeated Poles}{pfd_ex}
Consider
\begin{equation}
R(z) = \frac{4z + 4}{z(z - 1)(z - 2)^2} = \frac{A^{(1)}}{z} + \frac{A^{(2)}}{z - 1} + \frac{A_0^{(3)}}{(z - 2)^2} + \frac{A_1^{(3)}}{z - 2}.
\end{equation}
\begin{itemize}
    \item Simple pole at $z = 0$:
    \begin{equation}
    A^{(1)} = \lim_{z \to 0} z R(z) = \frac{4}{(0 - 1)(0 - 2)^2} = -1.
    \end{equation}
    \item Simple pole at $z = 1$:
    \begin{equation}
    A^{(2)} = \lim_{z \to 1} (z - 1) R(z) = \frac{8}{1(1 - 2)^2} = 8.
    \end{equation}
    \item Double pole at $z = 2$:
    \begin{align}
    A_0^{(3)} &= \lim_{z \to 2} (z - 2)^2 R(z) = \lim_{z \to 2} \frac{4z + 4}{z(z - 1)} = \frac{12}{2(1)} = 6, \\
    A_1^{(3)} &= \lim_{z \to 2} \frac{d}{dz}\left[\frac{4z + 4}{z^2 - z}\right] = \lim_{z \to 2} \frac{4(z^2 - z) - (4z + 4)(2z - 1)}{(z^2 - z)^2} = \frac{4(2) - (12)(3)}{2^2} = -7.
    \end{align}
\end{itemize}
Hence:
\begin{equation}
R(z) = -\frac{1}{z} + \frac{8}{z - 1} + \frac{6}{(z - 2)^2} - \frac{7}{z - 2}.
\end{equation}
\end{example}


\section{The Exponential, Trigonometric, and Hyperbolic Functions}

\subsection{The Complex Exponential}

\begin{proposition}{Properties of the Complex Exponential Function}{exp_func_props}
For $z = x + iy$:
\begin{equation}
e^z = e^x(\cos y + i \sin y).
\end{equation}
\begin{itemize}
    \item $e^z$ is \textbf{entire} and satisfies $\frac{d}{dz} e^z = e^z$.
    \item Polar coordinates: $|e^z| = e^x > 0$ (so $e^z \neq 0$ for all $z \in \mathbb{C}$), and $\arg(e^z) = y + 2k\pi$.
    \item \textbf{Periodicity}: $e^z$ is periodic with pure imaginary period $2\pi i$:
    \begin{equation}
    e^{z_1} = e^{z_2} \iff z_1 = z_2 + 2k\pi i \quad (k \in \mathbb{Z}).
    \end{equation}
    In particular, $e^z = 1 \iff z = 2k\pi i$.
    \item \textbf{Fundamental Region}: $e^z$ is one-to-one on any horizontal strip of width $2\pi$:
    \begin{equation}
    S_k := \{z \in \mathbb{C} : (2k - 1)\pi < \Im(z) \le (2k + 1)\pi\}, \quad k \in \mathbb{Z}.
    \end{equation}
\end{itemize}
\end{proposition}

\subsection{Complex Trigonometric Functions}

\begin{definition}{Complex Sine and Cosine}{trig_defs}
For any complex number $z \in \mathbb{C}$, the sine and cosine functions are defined by:
\begin{equation}
\sin z := \frac{e^{iz} - e^{-iz}}{2i}, \qquad \cos z := \frac{e^{iz} + e^{-iz}}{2}.
\end{equation}
\end{definition}

\begin{proposition}{Properties of Sine and Cosine}{trig_props}
\begin{itemize}
    \item Both $\sin z$ and $\cos z$ are \textbf{entire functions}.
    \item \textbf{Derivatives}:
    \begin{equation}
    \frac{d}{dz} \sin z = \cos z, \qquad \frac{d}{dz} \cos z = -\sin z.
    \end{equation}
    \item \textbf{Fundamental Identities} (identical to the real case):
    \begin{align}
    \sin^2 z + \cos^2 z &= 1, \\
    \sin(z + 2\pi) = \sin z, &\qquad \cos(z + 2\pi) = \cos z, \\
    \sin(-z) = -\sin z, &\qquad \cos(-z) = \cos z, \\
    \sin(z_1 \pm z_2) &= \sin z_1 \cos z_2 \pm \cos z_1 \sin z_2, \\
    \cos(z_1 \pm z_2) &= \cos z_1 \cos z_2 \mp \sin z_1 \sin z_2.
    \end{align}
    \item \textbf{Zeros}: All zeros of $\sin z$ and $\cos z$ are strictly \textbf{real}:
    \begin{align}
    \sin z = 0 &\iff z = k\pi \quad (k \in \mathbb{Z}), \\
    \cos z = 0 &\iff z = \frac{\pi}{2} + k\pi \quad (k \in \mathbb{Z}).
    \end{align}
    \item \textbf{Unboundedness}: Unlike the real cosine, $\cos z$ is unbounded in $\mathbb{C}$. For $z = iy$:
    \begin{equation}
    |\cos(iy)| = \left|\frac{e^{-y} + e^y}{2}\right| = \cosh y \xrightarrow{y \to \pm\infty} +\infty.
    \end{equation}
\end{itemize}
\end{proposition}

\begin{definition}{Tangent, Cotangent, Secant, and Cosecant}{other_trig}
\begin{equation}
\tan z := \frac{\sin z}{\cos z}, \quad \cot z := \frac{\cos z}{\sin z}, \quad \sec z := \frac{1}{\cos z}, \quad \csc z := \frac{1}{\sin z}.
\end{equation}
Their domains of analyticity and derivatives are:
\begin{align}
\frac{d}{dz}\tan z &= \sec^2 z \quad \left(z \neq \frac{\pi}{2} + k\pi\right), & \frac{d}{dz}\cot z &= -\csc^2 z \quad (z \neq k\pi), \\
\frac{d}{dz}\sec z &= \sec z \tan z \quad \left(z \neq \frac{\pi}{2} + k\pi\right), & \frac{d}{dz}\csc z &= -\csc z \cot z \quad (z \neq k\pi).
\end{align}
\end{definition}

\subsection{Complex Hyperbolic Functions}

\begin{definition}{Complex Hyperbolic Functions}{hyperbolic_defs}
For any $z \in \mathbb{C}$:
\begin{equation}
\sinh z := \frac{e^z - e^{-z}}{2}, \qquad \cosh z := \frac{e^z + e^{-z}}{2}.
\end{equation}
They are entire functions satisfying:
\begin{equation}
\frac{d}{dz}\sinh z = \cosh z, \qquad \frac{d}{dz}\cosh z = \sinh z.
\end{equation}
\end{definition}

\begin{proposition}{Relations between Trigonometric and Hyperbolic Functions}{trig_hyperbolic_rel}
\begin{align}
\sin(iz) &= i \sinh z, & \cos(iz) &= \cosh z, \\
\sinh(iz) &= i \sin z, & \cosh(iz) &= \cos z.
\end{align}
Consequently:
\begin{equation}
\cosh^2 z - \sinh^2 z = 1.
\end{equation}
The remaining hyperbolic functions are:
\begin{equation}
\tanh z := \frac{\sinh z}{\cosh z}, \quad \coth z := \frac{\cosh z}{\sinh z}, \quad \operatorname{sech} z := \frac{1}{\cosh z}, \quad \operatorname{csch} z := \frac{1}{\sinh z}.
\end{equation}
\end{proposition}


\section{The Complex Logarithmic Function}

\subsection{Definition of the Multi-valued Logarithm}

\begin{definition}{The Logarithmic Function}{log_def}
For any $z \neq 0$, the \textbf{complex logarithm} $\log z$ is defined as the multi-valued function:
\begin{equation}
\log z := \log|z| + i \arg(z) = \log|z| + i\Arg(z) + 2k\pi i \quad (k \in \mathbb{Z}),
\end{equation}
where $\log|z|$ is the ordinary real natural logarithm of the positive real number $|z|$.
\end{definition}

\begin{example}{Values of $\log z$}{log_vals}
\begin{itemize}
    \item $\log 3 = \log 3 + 2k\pi i$.
    \item $\log(-1) = \log 1 + i(\pi + 2k\pi) = (2k + 1)\pi i$.
    \item $\log(1 + i) = \log\sqrt{2} + i\left(\frac{\pi}{4} + 2k\pi\right) = \frac{1}{2}\log 2 + i\left(\frac{\pi}{4} + 2k\pi\right)$.
\end{itemize}
\end{example}

\begin{proposition}{Algebraic Properties of Multi-valued $\log z$}{log_properties}
For non-zero complex numbers:
\begin{enumerate}
    \item $e^{\log z} = z$ (holds for every value of $\log z$).
    \item $\log(e^z) = z + 2k\pi i$ ($k \in \mathbb{Z}$).
    \item \textbf{Addition and Subtraction} (understood in the sense of set equality):
    \begin{align}
    \log(z_1 z_2) &= \log z_1 + \log z_2, \\
    \log\left(\frac{z_1}{z_2}\right) &= \log z_1 - \log z_2.
    \end{align}
\end{enumerate}
\end{proposition}

\subsection{The Principal Value \texorpdfstring{$\operatorname{Log} z$}{Log z} and Branches}

\begin{definition}{Principal Value of the Logarithm}{principal_log}
The \textbf{principal value} of the logarithm, denoted $\operatorname{Log} z$, is obtained by selecting the principal argument $\Arg(z) \in (-\pi, \pi]$:
\begin{equation}
\operatorname{Log} z := \log|z| + i \Arg(z).
\end{equation}
$\operatorname{Log} z$ is a single-valued function. Its line of discontinuity is the non-positive real axis $(-\infty, 0]$, called the \textbf{branch cut}.
\end{definition}

\begin{theorem}{Analyticity and Derivative of $\operatorname{Log} z$}{log_analytic}
$\operatorname{Log} z$ is analytic in the cut plane (or slit domain):
\begin{equation}
D^* := \mathbb{C} \setminus \{x \in \mathbb{R} : x \le 0\}.
\end{equation}
Its derivative in $D^*$ is:
\begin{equation}
\frac{d}{dz} \operatorname{Log} z = \frac{1}{z} \quad \forall z \in D^*.
\end{equation}
\end{theorem}

\begin{remark}[Failure of Identities for Principal Values]
Identities valid for multi-valued $\log z$ do not necessarily hold for the principal value $\operatorname{Log} z$.
For example, for $z_1 = z_2 = -1$:
\begin{equation}
\operatorname{Log}((-1)(-1)) = \operatorname{Log}(1) = 0, \quad \text{but} \quad \operatorname{Log}(-1) + \operatorname{Log}(-1) = i\pi + i\pi = 2\pi i \neq 0.
\end{equation}
\end{remark}

\begin{definition}{Branches of the Logarithm}{branches_log}
A \textbf{branch} of $\log z$ is any single-valued function $L(z)$, analytic in a domain $D$, such that $e^{L(z)} = z$ for all $z \in D$.
For any real angle $\tau$, the branch taking values in $(\tau, \tau + 2\pi]$ is:
\begin{equation}
L_\tau(z) = \log|z| + i \arg_\tau(z),
\end{equation}
which is analytic in the domain obtained by deleting the ray $\arg z = \tau$ from $\mathbb{C}$. In this domain:
\begin{equation}
\frac{d}{dz} L_\tau(z) = \frac{1}{z}.
\end{equation}
\end{definition}


\section{Complex Powers}

\begin{definition}{Complex Power $z^\alpha$}{complex_powers}
Let $\alpha \in \mathbb{C}$ and $z \neq 0$. The \textbf{complex power} $z^\alpha$ is defined by:
\begin{equation}
z^\alpha := e^{\alpha \log z} = e^{\alpha(\log|z| + i\Arg(z) + 2k\pi i)} \quad (k \in \mathbb{Z}).
\end{equation}
\end{definition}

\begin{proposition}{Multiplicity of Values of $z^\alpha$}{power_multiplicity}
The number of distinct values assumed by $z^\alpha$ depends on the exponent $\alpha$:
\begin{enumerate}
    \item \textbf{$\alpha \in \mathbb{Z}$ (an integer)}: $z^\alpha$ is \textbf{single-valued} (giving the standard power $z^n$).
    \item \textbf{$\alpha = \frac{p}{q} \in \mathbb{Q}$ (a reduced rational number, $q > 0$)}: $z^\alpha$ has \textbf{exactly $q$ distinct values}:
    \begin{equation}
    z^{p/q} = \sqrt[q]{|z|^p} \exp\left[i \frac{p}{q}(\Arg(z) + 2k\pi)\right], \quad k = 0, 1, \dots, q-1.
    \end{equation}
    \item \textbf{$\alpha$ is irrational or non-real complex}: $z^\alpha$ has \textbf{infinitely many distinct values}.
\end{enumerate}
\end{proposition}

\begin{example}{Infinitely Many Values of $(-2)^i$}{power_ex}
\begin{equation}
(-2)^i = e^{i \log(-2)} = e^{i(\log 2 + i(2k+1)\pi)} = e^{-(2k+1)\pi} e^{i \log 2} \quad (k \in \mathbb{Z}).
\end{equation}
\end{example}

\begin{definition}{Principal Branch and Derivative of $z^\alpha$}{principal_power}
The \textbf{principal branch} of $z^\alpha$ is defined using the principal value of the logarithm:
\begin{equation}
z^\alpha := e^{\alpha \operatorname{Log} z}.
\end{equation}
It is analytic in the slit domain $D^* = \mathbb{C} \setminus (-\infty, 0]$, with derivative:
\begin{equation}
\frac{d}{dz} (z^\alpha) = \alpha z^{\alpha - 1}.
\end{equation}
\end{definition}


\section{Inverse Trigonometric and Hyperbolic Functions}

\begin{proposition}{Inverse Trigonometric Functions as Logarithms}{inv_trig}
The inverse trigonometric functions are multi-valued functions expressible in terms of the complex logarithm:
\begin{align}
\sin^{-1} z &= -i \log\left[iz + (1 - z^2)^{1/2}\right], \\
\cos^{-1} z &= -i \log\left[z + (z^2 - 1)^{1/2}\right], \\
\tan^{-1} z &= \frac{i}{2} \log\left(\frac{i + z}{i - z}\right) = \frac{i}{2} \log\left(\frac{1 - iz}{1 + iz}\right) \quad (z \neq \pm i).
\end{align}
Their derivatives (for any consistently chosen branch) are:
\begin{align}
\frac{d}{dz}\sin^{-1} z &= \frac{1}{(1 - z^2)^{1/2}} \quad (z \neq \pm 1), \\
\frac{d}{dz}\cos^{-1} z &= -\frac{1}{(1 - z^2)^{1/2}} \quad (z \neq \pm 1), \\
\frac{d}{dz}\tan^{-1} z &= \frac{1}{1 + z^2} \quad (z \neq \pm i).
\end{align}
\end{proposition}

\begin{proposition}{Inverse Hyperbolic Functions as Logarithms}{inv_hyperbolic}
\begin{align}
\sinh^{-1} z &= \log\left[z + (z^2 + 1)^{1/2}\right], \\
\cosh^{-1} z &= \log\left[z + (z^2 - 1)^{1/2}\right], \\
\tanh^{-1} z &= \frac{1}{2} \log\left(\frac{1 + z}{1 - z}\right) \quad (z \neq \pm 1).
\end{align}
\end{proposition}
